\section{Discussion}
\label{sec:discussion}

\paragraph{Information, not capacity.}
The same pattern recurs at every level of the project. A 1.2-million-parameter
roster network over event embeddings finished level with a constant; a linear
probe on the same embeddings did no better; two team-form features beat both.
Tree ensembles on engineered game features did not beat Elo. A small residual
network that reads the dressed lineup did beat Elo, confirmed on eight seasons,
and a far more structured engine that models every skater's deployment and
every team's scoring process reached the same skill and no more. What moved
the result was information that Elo does not use, namely who is dressed and
what those players have done, rather than any particular architecture. The
exploratory split in Section~\ref{sec:res-h} points the same way: \neurhlH{}'s
gain is near zero when neither team is missing regulars and largest when one
team is missing far more than the other. We
expect further gains to come from new information available before a game,
such as tracking data, line deployment and injury reports with dates, rather
than from larger models of the information already used.

\paragraph{What preregistration changed.}
The protocol altered this project's conclusions in concrete ways. It turned an
apparently significant tune-window result into an exploratory one and forced
a one-shot confirmation, which then passed. It withdrew a season-level claim
that rested on mismatched seasons. It retired a gate that would have selected
the more confounded rating system. It kept two confirmatory windows unspent
when pre-gates failed, including \neurhlG{}'s sealed seasons, which remain
available to a better candidate. It recorded, before the gate ran, that \neurhlG{}
was expected to fail its pre-gate, so the failure cannot be explained away
afterwards. None of these outcomes flatter the models; all of them make the
confirmed results more credible.

\paragraph{Audits over gates.}
Both leaks in the event simulator passed every realism gate, and the gate
implementation errors in Section~\ref{sec:failures} passed our own review
until the results were written up in full. Gates check what a model does with
its inputs; they do not check the inputs or the checkers. The two mechanisms
that caught these problems were a consumer that had to supply every input
itself and a direct causality test, and the act of restating every gate
against its declaration for publication.

\paragraph{Limitations.}
The \neurhlG{} gate seasons had been scored at game level by \neurhlH{} before \neurhlG{}
existed, so the gate result is evidence, not confirmation. The evaluation
covers one league and seasons 2008--2026, whose most recent years include a
change in how shots are recorded. Market data were excluded by design and no
historical odds exist in the repository, so we cannot say how close the
models come to the market. \neurhlH{}'s head selects its regularisation on training
loss, a weakness kept for the confirmation. Several claims about the season
layer and player seasons rest on development evidence only. Neural training
is not bit-reproducible on the hardware used, which we address by hashing
checkpoints and reporting CPU inference, but a retrained model will differ
slightly from the published one. Finally, in-season expected goals for
2026--27 depend on a third-party shot file being published on time.

\paragraph{Future work.}
The live season will provide the first clean test of \neurhlG{} and a second,
prospective test of \neurhlH{}. The sealed seasons 2025--2026 remain available for a
single confirmatory test of a future \neurhlG{} candidate under a new plan, which
could spend it with one declared change after the failed pre-gate, as the
current plan allows, or with a new model that brings new information.
Candidate information sources include player tracking data, available from
2022 but absent from most of the history, deployment models that predict line
combinations, and injury and scratch reports with reliable timestamps.

\section{Conclusion}
\label{sec:conclusion}

Under a protocol that fixed windows, gates and consequences in advance, a
neural player-level model conditioned on the dressed lineup beat a tuned Elo
system on eight held-out NHL seasons by 0.0048 nats per game, with a
confidence interval that excludes zero comfortably and a gain that comes from
resolution. A chained player-game model beat each skater's own history on all
four of its targets. A single-game simulation engine reached the same
win-probability skill as the confirmed model while producing a coherent box
score, but did not exceed it, so its sealed test was kept for a future
version. Most of what we tried did not work, and the record of those failures,
including our own errors, is part of the result. All forecasts for the
2026--27 season are public before their games and will be scored once, at the
end of the season.

\paragraph{Code and data availability.}
Code, preregistrations, search ledgers, committed predictions and the
acceptance tests that re-derive them are at
\url{https://github.com/ph05/neurhl}. Raw play-by-play, shift and tensor files
are not redistributed; the repository's fetchers rebuild them from the public
sources named in Section~\ref{sec:data}.
